\lesson{13}{May 16 2022 Mon (20:01:59)}{Proving Trigonometric Identities}

Up until now, we've studied trigonometric identites that are useful. But there
are many other identites that aren't particularly important but they exist and
they offer us an opportunity to learn another skill, which is proving
mathematical statements.

Let's try and proof the following identity:

\[ \sin (x) = \frac{\tan (x)}{\sec (x)} \].

\begin{proof}
  \label{prf:sin_equals_tan_over_sec}

  We can start with either side of the equation, but it's usually most
  sensible to start with the more complicated side since it will be easier to
  manipulate it. In this example, $\frac{\tan (x)}{\sec (x)}$ is more
  complicated than $\sin (x)$, so let's start with $\frac{\tan (x)}{\sec (x)}$ 
  and try to manipulate it until it looks like $\sin (x)$.

  \begin{align*}
    \frac{\tan (x)}{\sec (x)} &= \frac{\frac{\sin (x)}{\cos (x)}}{\frac{1}{\cos (x)}} \\
                              &= \frac{\sin (x)}{\cos (x)} \times \frac{\cos (x)}{1} \\
                              &= \sin (x)
  .\end{align*}

  This is a \textit{proof} of the identit $\sin (x) = \frac{\tan (x)}{\sec
  (x)}$ since we've shown that the right side is equivalent to the left side.
\end{proof}

Let's now try to proof $\frac{\cos (\theta)}{1 - \sin (\theta)} = \frac{1 +
\sin (\theta)}{\cos (\theta)}$. The proof of this identity requires a
relatively commonly employed \textit{"trick"} that is important to become
familiar with. Sometimes it is useful to use the \textbf{conjugate} of some
part (often the denominator) of the expression. The \textbf{conjugate} of the
expression $a + b$ is the expression $a - b$, and vice versa. By multiplying
the denominator by its conjugate, allows us to manipulate the left side that
is equivalent to the right side.

\begin{proof}
  \label{prf:cos_over_1_minus_sin_equals_1_plus_sin_over_cos}

  \begin{equation*}
    \begin{alignedat}{3}
      \frac{\cos (\theta)}{1 - \sin (\theta)} &= \frac{\cos (\theta)}{1 - \sin (\theta)} \times \frac{1 + \sin (\theta)}{1 + \sin (\theta)} \qquad&& \textrm{Using the conjugate of $1 - \sin (\theta)$} \\
                                              &= \frac{\cos (\theta)(1 + \sin (\theta))}{(1 - \sin (\theta))(1 + \sin (\theta))} \\
                                              &= \frac{\cos (\theta) (1 + \sin (\theta))}{1 - \sin (\theta) + \sin (\theta) = \sin^{2} (\theta)} \\
                                              &= \frac{\cos (\theta)(1 + \sin (\theta))}{1 - \sin^{2} (\theta)} \\
                                              &= \frac{\cos (\theta)(1 + \sin (\theta))}{\cos^{2} (\theta)} \qquad&& \textrm{Since $1 - \sin^{2} (\theta) = \cos^{2} (\theta)$} \\
                                              &= \frac{1 + \sin (\theta)}{\cos (\theta)} \\
                                              &= \frac{1 + \sin (\theta)}{\cos (\theta)} \\
    \end{alignedat}
  .\end{equation*}

  This is a proof that $\frac{\cos (\theta)}{1 - \sin (\theta)} = \frac{1 +
  \sin (\theta)}{\cos (\theta)}$ since we've shown that the left side is
  equivalent to the right side.
\end{proof}

\newpage
